\documentclass[11pt]{article}
\usepackage[margin=0.82in]{geometry}
\usepackage[T1]{fontenc}
\usepackage{lmodern,microtype}
\usepackage{xcolor,booktabs,longtable,tabularx,array}
\usepackage{enumitem,listings,amsmath,amssymb}
\usepackage{hyperref}
\hypersetup{colorlinks=true,linkcolor=blue!55!black,urlcolor=blue!55!black}
\setlength{\parindent}{0pt}\setlength{\parskip}{5pt}
\setlist{leftmargin=*,itemsep=2pt,topsep=3pt}
\lstset{basicstyle=\ttfamily\small,breaklines=true,columns=fullflexible,
  frame=single,showstringspaces=false,keepspaces=true}
\newcommand{\code}[1]{\texttt{#1}}
\newcommand{\observed}{\textbf{Observed: }}
\newcommand{\boundary}{\textbf{Claim boundary: }}
\title{Plain2Metta in Detail\\Purpose, Architecture, Evidence Pipeline,\\Worked Examples, and Next-Phase Roadmap}
\author{Expanded technical report for Benjamin Goertzel}
\date{Revised 19 August 2026}

\begin{document}
\maketitle

\begin{abstract}
Plain2Metta turns a constrained Plain specification into a chain of
inspectable, immutable artifacts, generated MeTTa and Python, bounded runtime
observations, and obligation-specific validation evidence.  Its central idea
is not that natural-language intent can be proved automatically.  It is that
every interpretive and executable step can be made explicit, independently
reviewable, hash-bound to its inputs, and invalidated when an upstream byte
changes.  This report explains the implemented system for a reader who knows
the basic ideas of Plain, MeTTa, and Lean~4 but has not followed Plain2Metta's
development.  It covers the two historical layers, the complete artifact
lineage, Stages 1--11 of semantic validation, the G0--G6 evidence vocabulary,
the Hypothesis/TLC/Z3/Lean/Hyperon/Python backends, the API and browser path,
and three exact worked examples.  It also explains the conceptual purpose of
the larger Plain2Metta pursuit, summarizes lessons from the prototype, gives
design guidance for four substantially more realistic benchmark tasks, and
compares a single-agent work plan with a six-agent parallel plan.  It
separates demonstrated behavior from broader research ambitions and records
acceptance and VM2 deployment evidence.
\end{abstract}

\tableofcontents
\newpage

\section{Introduction: purpose, design rationale, and lessons so far}

\subsection{The larger purpose of Plain2Metta}

The long-term purpose of Plain2Metta is not merely to save the labor of typing
MeTTa programs.  It is to investigate whether human-scale statements of
software intent can become the stable, inspectable center of a development
process in which machines elaborate, formalize, implement, test, and revise
software without losing the relationship between what was requested and what
was built.  Plain supplies a comparatively accessible language for stating
concepts, requirements, assumptions, tests, and risks.  MeTTa supplies a
compositional symbolic substrate in which those objects and their
relationships can be represented, queried, transformed, and eventually
reasoned about.  Grounded languages such as Python supply efficient ordinary
computation; model checkers and solvers explore bounded formal consequences;
Lean supplies kernel-checked proofs for the fragment precise enough to merit
them.  Plain2Metta is the controlled bridge among these roles.

That bridge is important for AI-assisted development.  A capable model can
produce plausible code from prose, but plausibility is not traceability.  A
second model can review the code, but agreement between models is not proof of
shared meaning.  Tests can pass while testing an accidental interpretation;
a theorem can be true while formalizing the wrong proposition.  If the
original intent, intermediate interpretation, validation plan, generated
program, and evidence are transient prompt context, it is difficult to locate
where an error entered or to know what a later edit invalidates.  The pursuit
therefore asks a more demanding question than ``can an LLM translate Plain to
MeTTa?'': can a mixed human--machine process make each consequential semantic
commitment explicit, reviewable, executable where appropriate, and
retractable when its premises change?

A successful answer could support a form of continuously revisable software
in which requirements are not dead prose beside the implementation.  They
would remain linked to executable examples, properties, formal obligations,
runtime monitors, and generated components.  Changes could trigger
semantically targeted regeneration and revalidation rather than an undirected
rewrite.  MeTTa is attractive in this picture not only as a target language
but as a medium for representing the evolving network of requirements,
interpretations, implementations, evidence, conflicts, and uncertainty.

\subsection{Why the design is artifact-centered and fail-closed}

The prototype deliberately decomposes translation into immutable artifacts
instead of treating a model response as the compiler.  This follows from
four conceptual observations.

First, meaning is version-relative.  Approval of one set of source bytes does
not approve a reworded document, even if a human considers the edit minor.
Content-addressed identities and transitive invalidation make that dependency
machine-checkable.  They are not administrative overhead; they are the
mechanism by which evidence remains about the object it purports to support.

Second, different tools establish different kinds of facts.  Hyperon can
execute MeTTa reductions; Python can perform grounded computation; Hypothesis
can search generated cases and shrink counterexamples; TLC can exhaust a
declared finite transition system; Z3 can establish satisfiability or
unsatisfiability of an encoding; Lean can kernel-check a theorem.  None of
these automatically establishes that the encoding captures the author's
intent.  Backend ownership and the G0--G6 vector preserve these distinctions
instead of compressing heterogeneous evidence into a misleading confidence
number.

Third, interpretation and validation are correlated failure sources.  If the
same generation step invents the meaning, writes the implementation, and
chooses the tests, it can create a self-consistent but wrong world.  The
reviewed-contract and independently authored validation-plan boundaries do
not guarantee independence, but they make role separation possible and leave
an auditable object at which a human or another model can disagree.

Fourth, honest incompleteness is valuable.  Real specifications contain
ambiguity, missing policy, environmental assumptions, and claims that cannot
be settled by the available tools.  An explicit typed hole or Unknown verdict
is more useful than synthesized certainty because it directs the next
question.  Fail-closed admission is therefore a semantic feature: unsupported
text is prevented from inheriting evidence produced for a nearby supported
profile.

\subsection{What the prototype has taught us}

The prototype has produced several practical lessons that should guide the
next phase.

\begin{enumerate}
\item \textbf{Provenance is implementable, but exactness propagates.}
Exact UTF-8 spans, canonical serialization, content hashes, immutable review
targets, and downstream invalidation can be made to work end to end.  Doing so
revealed that seemingly small boundary choices---newline handling, duplicate
keys, path normalization, current-version checks, and review revocation---are
part of semantic correctness rather than peripheral hygiene.

\item \textbf{A narrow vertical is more informative than broad mock
coverage.}  Running three exact examples through the real API, artifact
store, approved plan, formal/property tools, Hyperon and Python observations,
server-side verdict composition, and browser exposed integration errors that
component demonstrations hid.  The next benchmarks should again be complete
vertical slices, not isolated impressive backend examples.

\item \textbf{Evidence must be per obligation.}  A program may satisfy output
shape while losing trace identity, enforce a happy path while mishandling
retries, or prevent leakage while failing to report uncertainty.  One overall
pass badge obscures these distinctions.  Obligation-specific ancestry and
grades must remain visible as examples become larger.

\item \textbf{Formal machinery is useful only when its claim boundary is
shown.}  The prototype successfully invokes Hypothesis, TLC, Z3, and Lean with
pinned, attributed artifacts.  It also demonstrates how easy it would be to
overstate them.  A bounded TLC check is not an unbounded proof; a Z3 result is
about an encoding; a Lean theorem is about its exact statement and
assumptions.  The UI and reports should expose scope, formulas, theorem text,
and non-applicability.

\item \textbf{Mutation and adversarial tests reveal evidence-pipeline bugs.}
Source edits, stale plans, missing approvals, backend hash mismatches, timeouts,
malformed evidence, runtime disagreement, and false properties exercised the
trust boundaries better than more happy-path examples would have.  Future
benchmarks need semantic mutants designed before implementation so that
validation is evaluated rather than merely demonstrated.

\item \textbf{The present bottleneck is semantic breadth, not infrastructure
existence.}  The artifact schemas, bounded adapters, composer, API vertical,
and acceptance machinery now exist.  The live system remains narrow because
admission and profiles are curated around three exact texts and elementary
semantic shapes.  The next development phase should generalize only what
realistic benchmarks force it to generalize.

\item \textbf{Parallel development requires frozen seams.}  The historical
11-stage plan was appropriate for building trust boundaries in dependency
order.  Repeating it serially for every new case would waste the modularity it
created.  Once benchmark contracts and shared interfaces are frozen, case
fixtures, lowerings, oracles, mutants, and documentation can proceed in
parallel, with continuous integration guarding the common spine.
\end{enumerate}

These lessons motivate the appendices.  They do not yet supply finished
benchmark specifications.  Instead they define how those specifications
should be authored and how their implementation can be organized without
confusing concurrency with the removal of dependencies.

\section{Orientation: what problem Plain2Metta solves}

Plain is useful because a technical intent can be written in a form that is
more disciplined than ordinary prose while remaining readable by a domain
expert.  MeTTa is useful because relationships, rules, and symbolic
transformations can be represented and executed as atoms.  Lean~4 is useful
because a precisely stated proposition can be checked by a small proof
kernel.  The difficult part is the space \emph{between} these representations.
Nothing about translating an English sentence into a MeTTa rule guarantees
that the rule captures what its author meant, and nothing about successfully
executing a program makes its output semantically correct.

Plain2Metta addresses that gap as an evidence-management and compilation
pipeline.  It asks, in order:

\begin{enumerate}
\item What exact bytes did the author review?
\item Which concepts, requirements, tests, contracts, assumptions, and holes
      were extracted from those bytes?
\item Which interpretation was approved, and against which exact version?
\item What code was generated from that interpretation?
\item Under what runtime and resource bounds did that code execute?
\item Which independent examples, properties, models, solver obligations, or
      theorems support each requirement?
\item What remains unknown, conditional, or operationally unobserved?
\end{enumerate}

This framing matters.  A one-shot ``Plain in, MeTTa out'' demo tends to hide
interpretation choices.  Plain2Metta instead preserves those choices as
artifacts.  A reviewer can trace a verdict backward to evidence, from evidence
to an approved validation plan and implementation, and from those objects to
the exact source clauses that motivated them.

\subsection{What the current release does and does not claim}

The merged release implements a source-preserving, content-addressed,
fail-closed pipeline and a real evaluation UI vertical for three exact,
reviewed examples.  It includes a finite typed contract calculus, strict
semantic artifacts, bounded formal and runtime backends, server-side evidence
composition, and adversarial invalidation tests.

\boundary It is not a general natural-language program synthesizer and does
not prove that arbitrary Plain prose has been understood.  The three live
examples use exact reviewed profiles.  A reworded or unsupported document is
not allowed to inherit their evidence.  Passing examples are bounded evidence,
not a universal proof.  A Lean theorem proves only the theorem actually
encoded, under the assumptions actually admitted.  A successful deployment
canary is an operational observation over that canary interval, not permanent
production correctness.

\section{The two layers in the repository}

The public repository is named Plain2Metta; the internal package and
intermediate representation retain the name \code{specatom\_hs}.  Two layers
coexist because the project evolved from a conservative source compiler into
a larger reviewed artifact workflow.

\subsection{The earlier source-preserving compiler}

The \code{plain\_to\_metta} layer indexes a Plain file without throwing away
its origin.  It records file digests, sections, bullet items, concepts,
requirements, tests, and exact byte/line spans.  Stable identities and
provenance let a downstream atom point back to the characters that produced
it.  Conservative passes recognize explicit markers such as concepts,
requirements, coverage labels, claims, assumptions, evidence, risks,
mitigations, owners, deadlines, data-flow relations, temporal ordering, and
selected ML/security methodology cues.

The emphasis is conservative extraction.  Unsupported predicates, ambiguous
coverage, missing evidence, placeholder owners, vague deadlines, unresolved
references, or raw-text-only semantics become questions or refusals rather
than confident invented atoms.  The PeTTa/MeTTa backend emits only records
admitted by its profile; malformed arity, broken source links, unsupported
semantic levels, or unsafe scalar encodings fail closed.

This layer answers: ``What structure can we reproducibly preserve and expose
from the source?''  It intentionally does not answer: ``What executable
program should arbitrary prose mean?''

\subsection{The v2 immutable project and semantic-validation layer}

The \code{specatom\_hs} v2 layer turns translation into a staged project.  It
adds immutable artifact versions, SHA-256 content hashes, exact upstream
references, hash-bound review decisions, logical-IR admission, generated
output handoff, bounded test results, semantic artifacts, formal/runtime
evidence, traceability, and transitive invalidation.

The two layers are complementary.  The first preserves source structure and
supplies conservative semantic scaffolding.  The second controls what may be
promoted from a reviewed source into an interpretation, implementation, test,
or verdict.  Existing schema-v1 artifacts remain readable; new semantic
meaning is introduced through new versioned schemas rather than silently
changing an old envelope.

\section{The ordinary v2 artifact chain}

Before discussing semantic-validation Stages 1--11, it helps to understand
the underlying project chain.  A typical project contains the following
ordered artifacts:

\begin{center}
\small
\begin{tabularx}{0.96\textwidth}{@{}>{\bfseries}lX@{}}
\toprule
Artifact & Purpose \\
\midrule
Original specification & Exact submitted Plain bytes and their hash. \\
Elaborated specification & A proposed clarification/expansion bound to the original. \\
Test specification & Proposed tests and acceptance conditions, independently reviewable. \\
Reviewed snapshots & Exact approved elaborated and test bytes; comments and decisions bind to versions. \\
Logical IR & Typed declarations, contracts, obligations, dependencies, provenance, and explicit holes; no executable body. \\
Logical review & Machine-readable findings, reviewer dispositions, rationale, and compile-blocking state. \\
Compiler request/log & Exact reviewed inputs, guidance, provider attribution, response schema, and request digest. \\
Compiler output & Inert generated files with safe paths, hashes, sizes, and requirement traceability. \\
Sandbox handoff & Exact admitted generated bundle and bounded adapter request. \\
Test result & Per-test observation, status, output, duration, covered requirement IDs, adapter identity, and request digest. \\
Traceability report & Reconstruction from source through tests, classifying each requirement as passing, failing, skipped, or untested. \\
\bottomrule
\end{tabularx}
\end{center}

Each arrow is explicit.  An artifact reference contains an identity, version,
kind, and content hash.  Deserialization recomputes hashes and rejects forged
or malformed provenance.  Generic artifact creation cannot bypass the
special admission gates for logical IR, compiler output, or results.

\subsection{Review is attached to bytes, not filenames}

A mutable filename such as \code{requirements.plain} is not a trustworthy
review target: its contents can change after approval.  Plain2Metta attaches
annotations and decisions to an immutable artifact reference.  The Phase-3
review compares exact elaborated and test snapshots and records reviewer,
timestamp, target, decision, and rationale.  Logical findings likewise bind to
the exact current logical IR.  Repairing or waiving a critical finding creates
a new review version; it does not rewrite history.

Compilation is admitted only if the current reviewed inputs match the current
IR, the exact IR is approved, and no critical finding remains open or deferred.
If the source, elaboration, tests, approval, or IR changes, descendants become
stale.  This is \emph{transitive invalidation}: a verdict cannot remain current
when its semantic ancestry no longer exists.

\subsection{Generated code is inert until a separate handoff}

The compilation boundary stores generated text but marks it
\code{executed:false}.  It neither publishes files nor imports Python nor
evaluates MeTTa.  Paths must be safe, fields strict, and each generated file
must carry requirement traceability.  Execution requires a separate sandbox
handoff and an explicitly injected out-of-process adapter.  The coordinator
makes one call, supplies no provider selection or retry policy, and rechecks
the current ancestry after the call before admitting a result.

This separation limits authority.  A compiler may propose bytes; it does not
gain ambient authority to execute them.  An evidence runner may execute a
specific admitted request; it cannot reinterpret the source, approve a plan,
or mutate the project.

\section{The general semantic-validation model}

The semantic vertical extends the artifact spine with four key concepts.

\begin{description}
\item[SemanticContract] A typed, finite statement of inputs, output,
preconditions, postconditions, invariants, effects, temporal constraints,
assumptions, and unresolved holes, with source-clause ancestry.
\item[ValidationObligation] A particular claim linked to a contract and source
clauses, with required evidence grade, admissible methods, finite domain,
assumptions, severity, and unresolved state.
\item[ValidationPlan] Independently authored examples, generators,
properties, state models, differential oracles, and formal tasks, approved
against exact contract and source hashes.
\item[RuntimeEvidence and ValidationVerdict] Tool-attributed observations and
a conservative composition for one exact obligation, implementation, plan,
and set of evidence artifacts.
\end{description}

The contract language is deliberately finite and typed.  Its core supports
literals, variables, state reads, Boolean and numeric operators, bounded
quantification, trace predicates, explicit UTC time, approximation, and
approved assumption references.  Unknown operations, free prose, free
variables, malformed types, unapproved assumptions, and typed holes block
execution.  A hole is an honest record of unresolved meaning, not executable
placeholder code.

\subsection{Independence and correlated error}

The design separates interpretation author, implementation author, validation
author, and evidence runner.  The validation author may inspect approved
source and contracts but should not derive expected answers from generated
implementation bodies.  Otherwise the same misunderstanding can occur in
both implementation and oracle.  Full independence is sometimes impossible;
when dependencies are shared, they are named and the evidence claim is
limited.

The exact-profile evaluation fixtures are a useful example of this caution.
They provide deterministic and inspectable G3 oracles, but their generator
and expected outputs are curated in a shared profile data file.  That supports
the three reviewed examples; it is not evidence of general semantic induction.

\section{Evidence is a vector: G0 through G6}

Plain2Metta does not collapse unlike evidence into a single confidence number.
Each obligation receives a Boolean vector with exactly G0--G6 plus status,
assumptions, holes, counterexamples, attribution, and residual risk.

\begin{longtable}{@{}p{0.12\textwidth}p{0.80\textwidth}@{}}
\toprule
Grade & Meaning and common misinterpretation \\
\midrule
\endhead
G0 preserved & Exact reviewed source bytes, identities, hashes, and ancestry
are known.  This proves provenance, not meaning. \\
G1 structured & Contract and validation-plan representations pass schema,
typing, provenance, and consistency checks.  A well-formed interpretation may
still be the wrong interpretation. \\
G2 executable & An admitted artifact starts and terminates under its declared
bounded runtime.  Exit zero alone says nothing about expected behavior. \\
G3 examples & Independently specified concrete examples pass, including
complete normalized output where required.  Coverage is limited to those
examples. \\
G4 properties & Property-based, metamorphic, state-machine, differential, or
bounded model checks pass over a recorded domain, budget, and seed/model.
This remains bounded evidence. \\
G5 proved & A precisely stated formula is established by a named solver,
exhaustive finite checker, or proof kernel under explicit assumptions and
bounds.  It is not automatically a proof of the original prose. \\
G6 observed & A deployed system satisfies monitored obligations over a
recorded operational interval.  Observation does not imply universal future
behavior and is not stronger than proof in a simple linear sense. \\
\bottomrule
\end{longtable}

The vector form prevents silent promotion.  G2 does not imply G3; G4 does not
imply G5; and G6 is operationally different from G5.  A policy may require a
particular combination, but the response retains the components.  The legacy
\code{semantically\_validated} Boolean exists only for compatibility with the
old three-profile UI and denotes that its exact dual-runtime G3 oracle passed.

\section{Stages 1--11, step by step}

Stages 1--8 execute per supported evaluation request.  Stage 9 projects their
stored evidence to the API.  Stage 10 is immutable release calibration
metadata validated at service construction, not rerun per request.  Stage 11
is a clean-checkout release gate, not a runtime stage.

\subsection{Stage 1: preserve source and materialize semantic artifacts}

The service accepts only an exact bundled reviewed source.  Stage 1 creates
strict \code{plain2metta-semantic-artifact/v1} envelopes for contracts and
obligations, binding every identity and SHA-256 hash to ordered upstream
references and source clauses.  Canonical IDs cover complete documents.
Storage validates the envelope against immutable project ancestry rather than
trusting caller-supplied references.

Adversarial tests reject forged, partial, extended, duplicated, stale,
path-confused, hash-mismatched, or unknown-version documents.  A source-byte
mutation invalidates the contract $\rightarrow$ obligation $\rightarrow$ plan
$\rightarrow$ evidence $\rightarrow$ verdict chain.  Stage 1 therefore
establishes G0-style provenance and the substrate for G1; it does not interpret
arbitrary prose correctly by itself.

\subsection{Stage 2: finite contract calculus}

Stage 2 validates and lowers the typed contract calculus, runs its
deterministic reference interpreter, and can produce a canonical MeTTa
projection.  The interpreter defines what supported predicate forms mean,
rather than delegating semantics to whichever backend happens to run them.
Execution ownership is explicit: symbolic relations may belong to MeTTa;
ground numeric or systems operations may belong to Python.

A typed hole, unsupported operation, invalid type, unapproved assumption, or
unbounded construct blocks promotion.  This is important because it prevents
prose from being smuggled into a purportedly executable field.  Stage 2's
success means ``this approved finite representation has deterministic
semantics,'' not ``the representation perfectly captures human intent.''

\subsection{Stage 3: independently authored plan, review, and approval}

Stage 3 constructs exactly one validation plan for the approved contracts and
obligations.  The plan names concrete examples, input generators, properties,
models, solver tasks, proof tasks, finite budgets, and coverage claims.  It is
reviewed as its own immutable artifact.  Approval binds the exact plan,
contracts, obligations, and source hashes.

Missing, rejected, revoked, edited, stale, or mismatched approval prevents
execution.  Validation authors are not given generated implementation bodies
by default, reducing correlated-oracle error.  A plan-byte change invalidates
all results derived from it.  This stage is where the human or independent
author can reject a plausible but incorrect interpretation before expensive
or misleading evidence is generated.

\subsection{Stage 4: property and example checks with Hypothesis}

Stage 4 executes only Hypothesis tasks named in the approved plan.  The pinned
identity is Hypothesis 6.138.15.  A generator is itself reviewed: its domain,
distribution, budget, shrinker, and seed belong in the evidence.  A failure
retains the minimized counterexample.  A pass reports bounded cases; it is not
rephrased as proof over an infinite domain.

The evaluation contract bounds this adapter to 3 seconds wall time, 2 CPU
seconds, 128 MiB memory, a 1 MiB file limit, and one process.  Malformed or
misattributed output, a request-hash mismatch, timeout, or failing property is
conservative evidence, never a pass inferred from process exit alone.

\subsection{Stage 5: finite-state checking with TLC}

Stage 5 runs an approved finite-state model when the obligation has an
applicable temporal/state interpretation.  The pinned checker is TLC 1.7.4;
the recorded \code{tla2tools.jar} SHA-256 is
\code{936a262061c914694dfd669a543be24573c45d5aa0ff20a8b96b23d01e050e88}.
It runs under a project-local Temurin JRE 21.0.12+8.

The adapter is bounded to 8 seconds wall time, 5 CPU seconds, 3 GiB memory,
an 8 MiB file limit, and one worker.  A counterexample trace is preserved.  If
TLC is genuinely inapplicable, the result is a justified \code{Unknown}, not
an invented pass and not an error disguised as non-applicability.

\subsection{Stage 6: SMT checking with Z3}

Stage 6 checks an approved SMT obligation with pinned Z3 4.15.3.  It can
establish satisfiability/unsatisfiability or preserve a model/counterexample
for the exact encoded formula.  Bounds are 5 seconds wall time, 4 CPU seconds,
512 MiB memory, and a 4 MiB file limit.  The evidence schema identifies the
tool, request digest, formula ancestry, observations, and interpretation.

Z3 evidence is not relabelled as a Lean theorem.  A solver result remains
dependent on the encoding and solver identity.  Unknown, timeout, malformed
evidence, or disagreement with another runtime prevents unsupported grade
promotion.

\subsection{Stage 7: Lean~4 kernel checking}

Stage 7 checks an approved theorem using Lean 4.33.0 and Mathlib v4.33.0 at
commit \code{db584cd6d46c92f209a44c0f1c829460d327499d}.  The recorded Lean
archive SHA-256 is
\code{4b3fb03c29a1e0a253fb1d11f9bae3725f19a0dc6fc09b3ea16d2c9df3349e2c}.
It is bounded to 45 seconds wall time, 30 CPU seconds, 8 GiB memory, a 64 MiB
file limit, and one thread.

The important distinction is between kernel checking and prose validation.
The kernel can establish the encoded proposition under imports and
assumptions.  Plain2Metta retains the source/contract/theorem ancestry so a
reviewer can inspect whether that proposition addresses the intended
obligation.  Failed compilation, timeout, a theorem that does not close, or a
hash mismatch remains failure/unknown evidence.  TLC or Z3 output cannot be
presented as Lean evidence merely because it supports a similar claim.

\subsection{Runtime evidence: Hyperon MeTTa and bounded Python}

The Stage 4--7 orchestration also admits exact plan-authorized implementation
observations.  Generated MeTTa runs with pinned Hyperon CLI 0.2.10, under a
2-second timeout.  Generated Python runs in an isolated subprocess with a
2-second timeout, 2 CPU seconds, 128 MiB memory, a 1 MiB file limit, and one
process.  Neither adapter accepts a client-selected executable, shell command,
network capability, credential, or arbitrary host path.

The dual-runtime examples compare complete normalized stdout against reviewed
expected output.  Runtime success is G2; matching an independent example is
G3.  A zero exit code cannot override wrong output, and disagreement between
MeTTa and Python is a failure rather than something averaged away.

\subsection{Stage 8: conservative cross-tool composition}

Stage 8 joins only admitted evidence to each exact obligation.  It checks
kind, attribution, request hash, tool identity, ancestry, applicability,
observations, assumptions, and required method.  The composer produces status
(pass, fail, unknown, or blocked), the exact G0--G6 vector, counterexamples,
holes, assumptions, and residual risk.

No single passing backend can fill missing required evidence.  A property pass
cannot stand in for a proof; a Lean proof cannot fabricate operational
observation; an inapplicable TLC model cannot become a pass.  Failures and
unknowns remain visible.  This stage owns the grade decision; neither the API
route nor browser JavaScript may infer it.

\subsection{Stage 9: evidence projection for API and browser}

Stage 9 reloads the canonical chain and returns schema
\code{plain2metta-evaluation-evidence/v1}.  Its \code{ancestry.stages} reports
1 through 8.  Every node includes stage, kind, artifact ID, content hash,
state, and upstream references.  Verdicts bind obligation, implementation,
approved plan, evidence refs, and a complete Boolean G0--G6 vector.  The
projection also supplies plan/review/approval state, backend attribution,
counterexamples, assumptions, holes, residual risk, and Stage-10 metadata.

Artifact bodies are not dumped into the default response.  They remain behind
exact-ID, exact-hash, authorized download boundaries.  The browser renders the
server objects.  An independent audit found that it reads
\code{grade\_achieved.vector}; the earlier client-side grade inference and
hard-coded G4--G6 behavior are absent.

\subsection{Stage 10: calibrated release metadata}

Stage 10 is not recomputed for every request.  A strict checked-in data
artifact records release \code{plain2metta-r02-stage10-20260817}, corpus
schema/hash, acceptance examples, categorized mutations, and calibration
result.  Service construction validates the bytes once.  Missing, malformed,
changed, or unbound metadata makes the service unavailable rather than
selecting convenient defaults.

The fixed relevant mutation corpus scored 8/8 killed.  One cosmetic-wording
mutation was classified as non-semantic.  The score is evidence about this
finite corpus and policy; it is not a probability that the system is correct.

\subsection{Stage 11: clean-checkout release acceptance}

Stage 11 is an external acceptance gate.  Its script exercises the real
\code{/api/evaluate} route, Stage 1--9 traversal, Stage-10 consumption, pinned
backends, adversarial invalidation, browser framing, live examples,
compilation, and repository hygiene.  It also checks that no credential-like
content, oversized changed file, whitespace error, or unexpected untracked
artifact entered the candidate.

At accepted task head
\code{1755dddcb31dc02c04d7db35dec01ba1fb6b9215}, the focused real-route suite
passed 106/106 and the complete provider-free suite passed 777/777.  The fresh
clone rebuilt 3,015 pinned Lean jobs and passed.  A separate static audit
confirmed that the Flask route calls the Stage 1--8 orchestrator, Stage 9
contains Stage-10 metadata, and the browser consumes server grades.  These are
release-acceptance observations, not proof of unrestricted Plain semantics.

\section{The HTTP request and browser path}

The public evaluation surface is deliberately narrow.  \code{POST
/api/evaluate} accepts exactly a duplicate-free JSON object with one field:

\begin{lstlisting}[language={}]
{"text": "<UTF-8 Plain source>"}
\end{lstlisting}

Both request and decoded source are bounded to 128,000 bytes.  The text must
byte-match one of the three files returned by \code{GET /api/examples}.
Blank, mutated, reworded, malformed, duplicate-key, oversized, or unsupported
input fails closed without promoting partial state.  The client cannot supply
artifact IDs, plans, approvals, expected answers, backend choices, grades,
credentials, executable paths, or download authority.

For accepted input, the server creates an isolated canonical project chain,
executes the ordered vertical, reloads the stored evidence, and then runs the
compatibility/reference evaluator.  Repository failure, concurrent ancestry
change, invalid evidence, or any trust-boundary error aborts atomically.  There
is no retry, fallback provider, network access, host-shell authority, or
partial success committed as current.

The browser is therefore a presentation layer, not an evidence authority.  It
shows ancestry, reviews, tools, observations, counterexamples, assumptions,
holes, risks, and the G vector returned by the server.  This design makes a UI
bug less capable of upgrading evidence: the canonical verdict remains on the
server and in immutable artifacts.

\section{Worked example 1: greeting with trace preservation}

The first example is intentionally small.  It demonstrates that even a pure
transformation has two distinct requirements and that trace preservation is
not implied merely by producing greeting text.

\subsection{Exact Plain source}

\begin{lstlisting}[language={}]
***definitions***
- :Greeting: is a text response associated with a request trace.

***requirements***
- [id:GREET-1] The service shall return a :Greeting: for each request.
- [id:GREET-2] The :Greeting: shall retain the request trace identifier.
\end{lstlisting}

The exact curated profile is \code{greeting-trace-v1} with shape
\code{pure}.  The IDs matter: assertions link separately to GREET-1 and
GREET-2, so an output that returns ``hello'' but drops \code{trace-123} cannot
discharge both.

\subsection{Generated/reference MeTTa and expected observation}

\begin{lstlisting}[language={}]
(= (make-greeting $text $trace) (Greeting $text $trace))
!(make-greeting "hello" "trace-123")
\end{lstlisting}

The reviewed expected Hyperon output is:

\begin{lstlisting}[language={}]
[(Greeting "hello" "trace-123")]
\end{lstlisting}

The rule exposes the semantic shape directly: the two inputs become two
fields of a \code{Greeting} atom.  The query supplies a concrete text and
trace.  Exact-output comparison detects missing trace, changed spelling,
additional unexpected output, or a different constructor.

\subsection{Generated/reference Python and expected observation}

\begin{lstlisting}[language=Python]
def make_greeting(text, trace):
    return {'text': text, 'trace': trace}
def main():
    value = make_greeting('hello', 'trace-123')
    print(f"Greeting(text={value['text']},trace={value['trace']})")
\end{lstlisting}

Expected stdout is exactly:

\begin{lstlisting}[language={}]
Greeting(text=hello,trace=trace-123)
\end{lstlisting}

The assertions are ``response text equals hello'' for GREET-1 and ``response
retains trace-123'' for GREET-2.  The Python representation is not treated as
the definition of MeTTa semantics; it is a separately executed grounded
reference shape.  Agreement gives stronger concrete-example evidence than
one runtime alone, while the shared curated profile remains a documented
dependency.

\section{Worked example 2: task validation and owner preservation}

The second example introduces a conditional rejection path and a stateful
domain shape.  It separates three obligations: store valid data, reject
invalid data, and preserve ownership.

\subsection{Exact Plain source}

\begin{lstlisting}[language={}]
***definitions***
- :Task: is a named unit of work owned by a :User:.

***requirements***
- [id:TASK-1] A :User: shall add a valid :Task: to the task list.
- [id:TASK-2] The service shall reject a :Task: without a name.
- [id:TASK-3] The service shall preserve the owner of every stored :Task:.
\end{lstlisting}

The profile \code{task-validation-owner-v1} has shape \code{stateful}.  Its
examples exercise both the accepted and rejected branch rather than deriving
rejection behavior from the happy path.

\subsection{Generated/reference MeTTa and expected observations}

\begin{lstlisting}[language={}]
(= (add-task $name $owner)
   (if (== $name "")
       (Rejected "missing-name")
       (StoredTask $name $owner)))
!(add-task "Write report" "alice")
!(add-task "" "alice")
\end{lstlisting}

Expected output is two complete results:

\begin{lstlisting}[language={}]
[(StoredTask "Write report" "alice")]
[(Rejected "missing-name")]
\end{lstlisting}

The first result covers successful storage and shows that \code{alice} is
preserved.  The second is a negative boundary example.  Merely terminating or
returning some \code{Rejected} atom is not enough; the normalized output must
contain the reviewed \code{missing-name} reason.

\subsection{Generated/reference Python and expected observations}

\begin{lstlisting}[language=Python]
def add_task(name, owner):
    return ({'status':'rejected', 'reason':'missing-name'}
            if not name else
            {'status':'stored', 'name':name, 'owner':owner})
def main():
    valid = add_task('Write report', 'alice')
    invalid = add_task('', 'alice')
    print(f"stored={valid['name']} owner={valid['owner']}")
    print(f"rejected={invalid['reason']}")
\end{lstlisting}

Expected stdout is:

\begin{lstlisting}[language={}]
stored=Write report owner=alice
rejected=missing-name
\end{lstlisting}

The three linked assertions are: a valid named task is stored (TASK-1), a
blank name is rejected (TASK-2), and stored owner equals \code{alice}
(TASK-3).  This example also illustrates why grades are per obligation.  A
program could pass TASK-1 and TASK-2 while accidentally replacing every owner
with a default user; TASK-3 must remain independently visible.

\section{Worked example 3: forecast split, horizon, and baseline}

The third example addresses temporal lineage and an elementary anti-leakage
constraint.  It is not a forecasting benchmark.  Its purpose is to make data
split order, horizon, baseline identity, and reporting observable.

\subsection{Exact Plain source}

\begin{lstlisting}[language={}]
***definitions***
- :Observation: is a timestamped numeric measurement.
- :Forecast: is a prediction for a declared future horizon.

***requirements***
- [id:FORECAST-1] The pipeline shall train only on :Observation:
  records earlier than the evaluation window.
- [id:FORECAST-2] The pipeline shall emit one :Forecast: for the
  declared horizon.
- [id:FORECAST-3] The evaluation shall compare each :Forecast:
  with a named baseline.
- [id:FORECAST-4] The evaluation shall report the split timestamps
  used for each result.
\end{lstlisting}

The profile is \code{forecast-split-horizon-baseline-v1}, shape
\code{temporal-lineage}.

\subsection{Generated/reference MeTTa and expected observations}

\begin{lstlisting}[language={}]
(= (eligible-train $time $start) (< $time $start))
!(eligible-train 10 20)
!(eligible-train 30 20)
\end{lstlisting}

Expected Hyperon output is:

\begin{lstlisting}[language={}]
[True]
[False]
\end{lstlisting}

This isolates the ordering predicate behind FORECAST-1: time 10 is eligible
before evaluation start 20, while future time 30 is not.  The MeTTa fixture
does not pretend to train a model.  It validates the symbolic split predicate,
leaving numerical/reporting behavior to the grounded Python example.

\subsection{Generated/reference Python and expected observations}

\begin{lstlisting}[language=Python]
def main():
    train = [x for x in [10, 19, 30] if x < 20]
    print(f'train={train} evaluation_start=20')
    print('forecast=42 horizon=24h')
    print('baseline=seasonal-naive delta=2')
\end{lstlisting}

Expected stdout is:

\begin{lstlisting}[language={}]
train=[10, 19] evaluation_start=20
forecast=42 horizon=24h
baseline=seasonal-naive delta=2
\end{lstlisting}

The assertions are: future timestamp 30 is excluded (FORECAST-1), one 24-hour
forecast is emitted (FORECAST-2), the named \code{seasonal-naive} baseline is
compared (FORECAST-3), and split timestamp 20 is reported (FORECAST-4).
The constant forecast value 42 and delta 2 are fixture values, not a scientific
claim.  The useful result is that the critical lineage fields are explicit and
exactly checked.

This example shows backend ownership clearly.  MeTTa represents a temporal
eligibility relation.  Python grounds list filtering and formatted evaluation
output.  Both observations retain requirement links.  A future, richer
forecasting profile could add generators, uncertainty methods, real data
identifiers, and formal leakage obligations without changing the provenance
model.

\section{Adversarial and fail-closed behavior}

A trustworthy pipeline must be tested not only on its happy path but on ways
evidence could be accidentally or maliciously upgraded.  The accepted suite
covers the following classes.

\subsection{Source and approval attacks}

Exact source mutation, unsupported rewording, changed contracts or
obligations, changed plan bytes, stale immutable review, missing approval,
rejected approval, and approval revocation all invalidate the appropriate
descendants.  Tests check proportional invalidation: upstream state is
preserved, but evidence and verdicts that depended on it cease to be current.
This prevents ``approved once, valid forever'' behavior.

\subsection{Evidence substitution and framing attacks}

The system rejects malformed, partial, duplicate-key, unknown-version,
path-confused, hash-mismatched, or misattributed artifacts.  Backend results
must carry the exact request digest and tool identity.  A TLC/Z3 observation
cannot masquerade as Lean proof evidence.  Artifact-body download requires
the exact ID and hash plus authorization; metadata routes do not confer
mutation or execution authority.

\subsection{Behavioral and backend failures}

Wrong expected output, cross-runtime disagreement, a failing property,
failing TLC model, failing Lean theorem, backend timeout, malformed Stage-4
evidence, or resource exhaustion remains fail/unknown/blocked according to the
approved policy.  Process exit zero cannot override observations.  An
inapplicable formal method yields a justified unknown record.  Missing
required evidence prevents grade promotion even if another backend passes.

\subsection{Browser and route integrity}

Tests instrument the real Flask route rather than a detached service helper.
They assert ordered Stage 1--8 traversal, Stage-9 projection, Stage-10 metadata,
complete G0--G6 server vectors, duplicate-safe JSON, size bounds, exact source
admission, and absence of client-side grade synthesis.  This closes a common
integration gap in which well-tested backend code is silently bypassed by the
actual UI route.

\section{Acceptance evidence and deployment}

\subsection{Accepted and merged revisions}

The evaluation-UI task branch was accepted at
\code{1755dddcb31dc02c04d7db35dec01ba1fb6b9215}.  PR \#3 was then merged to
\code{main} as
\code{5ce102cb02e72f377314205b28f102e5fc39911f}.  The merge did not alter the
accepted head's content.

\observed The Stage-11 task-branch run passed 106/106 focused real-route tests
in 53.975 seconds and 777/777 complete provider-free tests in 70.501 seconds,
followed by byte compilation, diff, credential-pattern, large-file, and
untracked-file checks.  A fresh clone replay rebuilt 3,015 Lean jobs.  A
separate audit confirmed the intended server/browser code path.

These counts describe exact recorded environments and revisions.  Earlier
incremental runs include expected red gates and test-harness mistakes; they
were preserved rather than rewritten as successes.  Examples include a
missing \code{PYTHONPATH}, a test that patched a read-only \code{PosixPath}
method, and intermediate failures before the next stage existed.  Corrected
runs are separately recorded.  This makes the development evidence auditable
instead of presenting only a polished endpoint.

\subsection{VM2 deployment}

The merged commit is installed in an immutable release directory on the
existing VM2.  One enabled and active \code{plain2metta-vm2.service} worker is
bound to \code{127.0.0.1:8081}.  It is localhost-only, not publicly exposed.
No provider resource was provisioned, started, resized, or retained for this
deployment.

\observed Corrected live canaries passed \code{GET /}, health, examples, and
\code{POST /api/evaluate} for all three bundled sources.  Each response had
Stage 1--8 ancestry, Stage-9 schema
\code{plain2metta-evaluation-evidence/v1}, Stage-10 release
\code{plain2metta-r02-stage10-20260817}, and all server-derived G0--G6 keys.

Two prior deployment trials failed safely.  First, root-owned Lean cache trace
files were unreadable by the unprivileged service.  Second, shell command
substitution stripped the exact source's trailing newline, so source admission
correctly returned HTTP 422.  Each trial rolled back before correction.  The
newline incident is especially instructive: exact byte identity is not merely
documentation; it changed the route outcome as designed.

\section{Trust boundaries and security posture}

Trusted core operations are narrow: canonical serialization, hash
calculation, strict schema validation, exact-chain reconstruction, and
invalidation.  Plain prose, model/provider output, generated code, validation
plans, backend stdout/stderr, uploaded names, PR text, and runtime dependencies
are untrusted inputs.  Hyperon, Python, Hypothesis, TLC, Z3, Lean, Mathlib, and
their launchers are identified dependencies, not magically part of the
semantic trusted core.

The web client has no command authority.  Runtime adapters are injected and
command-free from the caller's perspective.  Network, credentials, arbitrary
host paths, and shell commands remain outside the request schema.  Bounded
JSON and duplicate-key rejection reduce framing ambiguity.  Atomic repository
admission prevents half-written evidence chains.  One-call/no-retry behavior
limits unreviewed mutation and cost; provider selection and credential policy
are deliberately not hidden in the core.

This does not make generated execution risk-free.  Resource limits, process
isolation, localhost binding, exact handoffs, and absence of ambient authority
are defense layers.  A production deployment would still need OS/container
hardening, monitoring, patch management, and a deliberate exposure policy.

\section{What has actually been accomplished}

The work can be summarized as a progression from syntax preservation to
evidence-bearing execution:

\begin{enumerate}
\item Exact Plain bytes, spans, stable identities, requirements, tests,
      concepts, semantic markers, questions, and conservative PeTTa atoms are
      preserved rather than flattened into opaque generated code.
\item Immutable project artifacts, SHA-256 provenance, strict persistence,
      version-bound comments/decisions, atomic writes, and transitive
      invalidation create a trustworthy state spine.
\item Reviewed elaboration/test snapshots and non-executable logical IR make
      interpretive decisions inspectable.  Critical gaps block compilation
      unless explicitly repaired or waived with attribution and rationale.
\item Provider-neutral one-call elaboration/compilation seams store inert
      output and exact provenance without bundling credentials, retries, or
      ambient execution.
\item Separate sandbox handoff/results and traceability reconstruct how each
      requirement maps to generated paths and passing/failing/skipped/untested
      tests.
\item Strict semantic contracts, obligations, validation plans, evidence,
      counterexamples, and verdicts provide a finite general framework beyond
      the original compatibility Boolean.
\item Hypothesis, TLC, Z3, Lean, Hyperon, and Python are integrated as
      distinct, pinned, bounded evidence sources.  Attribution and scope are
      retained rather than conflated.
\item The real evaluation API now traverses Stages 1--8, returns a Stage-9
      projection with Stage-10 calibration, and the browser renders the
      server's complete per-obligation G vector.
\item The adversarial matrix, clean-checkout Stage-11 gate, independent audit,
      merge, and localhost-only VM2 deployment have concrete recorded evidence.
\end{enumerate}

\section{Limitations and open research}

\subsection{Restricted semantic coverage}

The largest limitation is semantic breadth.  The finite calculus and explicit
profiles cover reviewed forms; they do not understand unrestricted English or
the full Plain language.  The earlier compiler intentionally extracts only
explicit markers and conservative patterns.  Rich policies, domain-specific
units, probabilistic claims, complex concurrency, open-world knowledge, and
underspecified operational behavior need new typed forms and review.

The three examples are pedagogically useful but tiny.  They demonstrate
provenance, negative cases, temporal split lineage, and dual-runtime evidence;
they do not demonstrate large programs, long-lived mutable state, distributed
systems, realistic ML training, or autonomous synthesis.

\subsection{Interpretation remains the central epistemic risk}

Formal machinery can verify the wrong formalization perfectly.  Hashes prove
identity, not adequacy.  Independent roles, source links, explicit holes, and
human approval reduce risk but do not eliminate it.  Future work should study
how reviewers compare source clauses to contracts, how disagreements are
represented, and how multiple independent interpretations can be evaluated
without mistaking consensus for truth.

\subsection{Evidence coverage and calibration}

The 8/8 mutation result is based on a fixed categorized corpus.  Larger and
less profile-shaped mutation sets are needed, including semantic-preserving
changes, subtle oracle corruption, backend attribution attacks, generator
bias, proof/contract mismatch, and operational drift.  Mutation score should
remain a diagnostic, not be converted into a correctness probability.

G6 requires real monitored obligations over an explicit interval.  A live
canary is useful deployment evidence, but sustained operational observation,
data retention, incident semantics, and privacy/security policies remain
deployment-specific work.

\subsection{Providers, UI breadth, and execution infrastructure}

The provider-neutral seams do not bundle a production model provider,
credential manager, retry/fallback policy, or cost policy.  The full v2 product
UI envisaged editors, workflow bars, logical-review drill-down, generated-code
inspection, results, and version sidebars; the shipped evaluation UI is the
narrow evidence vertical, not that full authoring environment.

The core deliberately lacks a general host/container executor.  The live
vertical uses explicitly configured bounded adapters.  Broader execution will
need a carefully reviewed sandbox and publication boundary.  Public exposure
of the VM2 service has not been authorized or implemented.

\subsection{Promising next investigations}

Useful next steps include: a richer library of typed contract patterns;
larger independently authored examples; explicit unit/dimensional semantics;
state-machine and temporal patterns with source-level explanations; proof
obligation generation whose adequacy can itself be reviewed; independent
validation-plan providers; differential execution across MeTTa runtimes;
metamorphic relations for symbolic transformations; realistic ML leakage and
uncertainty profiles; operational monitors that produce honest G6 records;
and usability studies of whether reviewers can find planted interpretation
errors from the evidence graph.

\section{How to read a Plain2Metta result responsibly}

For a practical evaluation, begin with the source and obligation, not the
headline.  Check that the bytes and clause IDs are the intended version.  Read
the contract, assumptions, and holes.  Verify who or what authored and
approved the validation plan.  Inspect which backend produced each evidence
item, its bounds, seed/model/theorem, and exact observation.  Treat a
counterexample as first-class.  Confirm that the verdict's plan and
implementation hashes are current.  Finally, interpret each G bit according
to its own meaning rather than imagining a scalar maturity score.

For the three bundled examples, the correct conclusion is specific: the exact
reviewed sources traversed the current immutable pipeline; their curated
MeTTa and Python fixtures produced exact expected observations under bounded
runtimes; the approved formal/property tasks and server composer produced the
recorded vectors; and the merged service replayed the path on VM2.  The
incorrect conclusion would be that Plain2Metta can now prove arbitrary Plain
documents equivalent to arbitrary MeTTa programs.

\section{Reproducibility record}

The public repository is \url{https://github.com/bgoertzel-sing/plain2metta}.
The accepted task head is
\code{1755dddcb31dc02c04d7db35dec01ba1fb6b9215}; merged \code{main} is
\code{5ce102cb02e72f377314205b28f102e5fc39911f}.  The main acceptance command
is \code{scripts/run-stage11-acceptance.sh} from a clean checkout with the
recorded project-local pinned dependencies and Lean prebuild.  The detailed
experiment records are retained in the project notebook under:

\begin{itemize}
\item \code{experiments/20260818T065419Z-plain2metta-evaluation-stage11-vertical-acceptance/};
\item \code{experiments/20260818T072517Z-plain2metta-evaluation-stage11-clean-checkout/};
\item \code{experiments/20260818T075514Z-plain2metta-evaluation-stage11-independent-audit/}; and
\item \code{experiments/20260818T153337Z-plain2metta-pr3-merge-vm2-deployment/}.
\end{itemize}

The source of truth for exact examples is
\code{examples/evaluation/*.plain}; generation/expected-output profiles are in
\code{src/specatom\_hs/evaluation\_profiles.json}; immutable Stage-10 data is
in \code{src/specatom\_hs/vertical\_acceptance.json}; and the route contract
is \code{docs/evaluation-vertical-contract.md}.  These files outrank this
expository narrative if a future revision changes behavior.

\section{Conclusion}

Plain2Metta's substantive result is an explicit chain of accountability from
Plain source to evidence.  Exact bytes and hashes prevent unnoticed drift;
reviewed contracts and plans expose interpretation; inert compilation and
separate bounded execution constrain authority; multiple backends retain
their distinct epistemic meanings; conservative composition prevents one
success from filling unrelated gaps; and the API/browser show the resulting
ancestry rather than a hand-waved badge.

The present system is deliberately narrow in semantic coverage but unusually
explicit about what it knows.  That is a sound foundation for broader
Plain-to-MeTTa research: new contract forms, generators, proofs, and operational
monitors can be added without abandoning provenance, role separation, exact
review, or fail-closed behavior.

\appendix

\section{Guidance for authoring the next realistic benchmark tasks}

This appendix is a design brief for the person or coding agent who will write
the benchmark specifications.  It intentionally stops short of supplying the
finished Plain documents, formal contracts, expected outputs, or complete
fixture data.  Those objects should be authored as separately reviewable
artifacts.  Prematurely filling them in here would make the implementation
appear more settled than it is and would weaken the separation between
benchmark intent and system construction.

\subsection{Rules that apply to every benchmark}

Each benchmark should describe a small but credible system rather than a
large fictional product.  It should be complicated enough that several
requirements interact and a superficially plausible implementation can be
wrong.  Its environment must be finite or fixture-bounded so that another
person can reproduce the result locally.  External services, private data,
and nondeterministic network dependencies should be represented by explicit
interfaces and recorded traces rather than required for the acceptance run.

The benchmark author should create the following objects before a coding
agent implements a new semantic profile:

\begin{enumerate}
\item a short domain narrative naming actors, state, boundaries, and the
      decision or service the software performs;
\item a Plain source document with stable requirement IDs, definitions,
      explicit assumptions, non-goals, and unresolved questions;
\item an obligation table stating what would count as Pass, Fail, Unknown,
      Blocked, and inapplicable for each requirement;
\item a proposed typed-contract representation, with unsupported meaning
      retained as typed holes rather than encoded by approximation without
      review;
\item a backend-ownership table: which claims are checked by reference
      interpretation, property tests, model checking, SMT, Lean, MeTTa,
      grounded execution, or operational observation;
\item deterministic fixtures and independent oracles, including exact
      expected observations where exactness is meaningful and tolerance or
      statistical procedures where it is not;
\item a preregistered mutant set containing realistic semantic errors, not
      only syntax corruption;
\item expected G0--G6 vectors per obligation, including justified zeros and
      acceptable Unknowns; and
\item an end-to-end acceptance recipe that reconstructs every result from
      exact source, tool versions, seeds, scopes, and artifact hashes.
\end{enumerate}

At least one requirement in each benchmark should be deliberately outside the
initial supported fragment.  The desired outcome for that requirement is an
honest question, typed hole, or Unknown verdict.  This tests whether broader
inputs cause the system to invent semantics.  At least one mutant should be
self-consistent across implementation and naive tests, so that it is caught
only by an independently stated invariant or oracle.

\subsection{Benchmark A: transactional reservation and payment saga}

\subsubsection*{Why this case}

This case should exercise state, idempotency, partial failure, compensation,
and externally visible effects.  A bounded reservation/payment workflow is
realistic enough to expose errors common in distributed services while still
admitting a finite model.  It extends the toy idempotency examples into a
system in which two resources must not silently diverge.

\subsubsection*{System envelope to specify}

The benchmark author should choose a concrete small workflow, such as
reserving one inventory item and authorizing one payment for a request bearing
an idempotency key.  Define the state machines for reservation, payment, and
the coordinating request.  State exactly which effects are durable, which
messages may be duplicated, delayed, reordered, or lost, where a crash may
occur, and what retry policy is assumed.  Decide whether compensation means
release, refund, cancellation request, or eventual manual review; these are
not interchangeable.

Candidate obligations should include: no completed order without both valid
effects; at most one captured payment per idempotency key; no oversold item;
duplicate requests return a result consistent with the first committed
outcome; recovery after each enumerated crash point; and an explicit terminal
or escalated state rather than silent limbo.  Liveness claims must state their
fairness and availability assumptions.  If the first benchmark only checks
safety in a finite scope, progress should remain Unknown rather than being
inferred from absence of a counterexample.

\subsubsection*{Expected representations and evidence}

The Plain document should identify actors, commands, events, states, effects,
and invariants with stable IDs.  The logical IR should preserve the distinction
between a requested, authorized, captured, released, and refunded effect.  Its
contract calculus should use bounded traces, explicit state transitions,
idempotency-key equality, and a finite inventory count.  MeTTa should represent
the workflow relations and permit queries such as whether a trace contains a
terminally consistent pair of effects.  Grounded Python should execute a
deterministic reference saga over scheduled event traces.

TLC is the natural owner of bounded interleavings, duplicate delivery, crash
points, recovery, and safety invariants.  Hypothesis should generate command
and failure schedules and shrink an observed violation.  Z3 can check local
state/invariant consistency and boundary conditions.  Lean should initially
be reserved for a small, reviewed theorem---for example, that a particular
pure transition preserves an at-most-once token invariant---rather than a
claim about the whole distributed workflow.  G6 should remain false unless a
separate deployed monitor observes real transactions over a declared interval.

\subsubsection*{Mutants and questions the author must settle}

Include mutants that: record an idempotency key after rather than before an
effect; capture twice after a retry; decrement inventory twice; acknowledge
success before durable commit; lose a compensation request; conflate payment
authorization with capture; recover from stale state; or reuse a key across
tenants.  Require source-linked counterexample traces.

Before finalizing the benchmark, settle the meaning of ``exactly once,'' the
atomicity available inside each service, the fate of permanently failing
compensation, the scope of tenant identity, and whether eventual progress is
part of the contract.  These questions should not be decided implicitly by a
generated model.

\subsection{Benchmark B: concurrent authorization with revocation}

\subsubsection*{Why this case}

This benchmark should test policy, identity, temporal ordering, cache
freshness, and race conditions.  Authorization looks like a Boolean function
in a toy example, yet real decisions depend on principal, resource, action,
delegation, policy version, context, and revocation time.  A stale allow is a
useful example of a program that can pass ordinary unit tests while violating
the intended security invariant.

\subsubsection*{System envelope to specify}

Define a small multi-tenant document or workflow service.  Include direct
roles, one bounded form of delegation, policy-versioned authorization
decisions, a cache, audit events, and revocation.  Specify whether a request is
judged at arrival, decision, or effect-commit time; how clock/order is modeled;
what consistency the policy store provides; and the maximum permitted cache
staleness, if any.

Candidate obligations should cover tenant isolation; default deny for missing
or malformed identity; action/resource-specific permission; delegation scope;
revocation behavior; cache invalidation; preservation of request trace and
policy version in the audit record; and non-disclosure of sensitive policy
details in a denial.  Include a policy conflict whose desired result is
explicitly unresolved until a precedence rule is approved.

\subsubsection*{Expected representations and evidence}

The logical IR should distinguish facts, derived permissions, denials,
delegations, policy versions, and effect events.  MeTTa is well suited to
representing and querying role/delegation relationships, provided negation and
closed-world assumptions are explicit.  Python can ground request processing,
cache behavior, and exact audit output.  Hypothesis should generate principals,
resources, action combinations, policy changes, and short concurrent traces.

TLC should explore orderings among authorization, revocation, cache refresh,
and effect commit.  Z3 can check finite policy consistency, unreachable allow
conditions, and tenant separation in the supported fragment.  Lean may prove
a small non-escalation or tenant-separation theorem for a precisely defined
policy evaluator.  The benchmark must show the exact theorem and its treatment
of stale external state; it must not describe the theorem as proving the
deployed authorization system secure.

\subsubsection*{Mutants and questions the author must settle}

Mutants should include: omit tenant from a cache key; treat missing identity as
anonymous read; allow role names to collide across tenants; apply revocation
only after cache expiry; widen delegated action scope; log the current rather
than decision-time policy version; perform the protected effect before the
final decision; and leak denial rationale.  Include a reworded Plain policy
that is plausible but unsupported and must fail admission rather than borrow
the curated profile.

The benchmark author must explicitly resolve denial precedence, decision-time
semantics, clock assumptions, cache bounds, delegation depth, audit privacy,
and whether revocation affects already-started operations.  A human reviewer
should be able to locate each choice in the source and resulting contract.

\subsection{Benchmark C: temporal ML evaluation and deployment decision}

\subsubsection*{Why this case}

This benchmark extends the existing temporal-lineage fixture into a realistic
data/estimation task.  It should test whether Plain2Metta can preserve dataset
identity, temporal splits, feature availability, baselines, uncertainty,
selection criteria, and deployment conditions without pretending that
symbolic execution itself validates a predictive model.

\subsubsection*{System envelope to specify}

Choose a small public or generated time series with an immutable local
snapshot.  Define an as-of forecasting or risk-estimation task, prediction
horizon, feature timestamps, train/validation/test windows, a named naive
baseline, a metric, an uncertainty procedure, and a decision rule.  Record
which transformations are fit using which window.  If model selection occurs,
separate selection data from final evaluation data.

Candidate obligations should include: no future or unavailable feature enters
training or inference; split timestamps and data hash are reported; all
competitors use identical evaluation cases; the named baseline is computed;
metric direction and aggregation are correct; uncertainty is estimated by a
declared method; missingness and degenerate windows fail explicitly; and a
deployment recommendation is conditional on preregistered thresholds rather
than narrated after seeing results.  Claims about future distribution shift
should normally remain assumptions or operational obligations.

\subsubsection*{Expected representations and evidence}

The IR should represent dataset versions, event time, availability time,
feature lineage, transforms, windows, horizons, metrics, comparators,
uncertainty procedures, seeds, and decision thresholds.  MeTTa should expose
lineage and eligibility relations and support queries about whether a feature
was available as of a forecast origin.  Grounded Python should perform the
actual numerical fitting and evaluation in an isolated, deterministic local
environment.

Hypothesis should generate boundary timestamps, missing values, duplicated
records, empty windows, and transformations to test leakage and metric
invariants.  Z3 may check finite ordering constraints and mutually inconsistent
window declarations.  Lean can prove a small property of split eligibility or
aggregation over a simplified data type.  TLC is appropriate only if a
stateful streaming/deployment protocol is included; otherwise it should be
marked inapplicable.  Scientific performance evidence is not automatically
G6: operational grade requires a monitored deployed interval and drift/error
policy.

\subsubsection*{Mutants and questions the author must settle}

Mutants should include: fit normalization on all data; leak the target through
a lag error; reverse the metric comparison; compare models on different rows;
silently drop hard cases; select and evaluate on the same window; misalign the
horizon; report a confidence interval computed over inappropriate independent
units; or change the data while retaining an old digest in the report.

The author must settle event versus ingestion time, timezone and boundary
conventions, missing-data policy, stochastic seed policy, resampling unit,
multiple-comparison treatment, practical versus statistical significance,
and what evidence would justify deployment.  The benchmark should reward an
honest inconclusive result when the fixture does not support a superiority
claim.

\subsection{Benchmark D: evidence-based policy decision under conflict and ignorance}

\subsubsection*{Why this case}

The first three cases emphasize operational correctness.  This case should
exercise MeTTa's prospective role in explicit reasoning with incomplete,
conflicting, differently sourced evidence.  It tests whether the pipeline can
preserve epistemic distinctions instead of collapsing absent evidence,
negative evidence, disagreement, and low confidence into one Boolean.

\subsubsection*{System envelope to specify}

Define a bounded decision-support setting such as approving an automated
action, escalating it to a human, or declining to decide.  Inputs should
include claims from sources with declared identity and provenance, positive
and negative observations, source dependence, policy constraints, and one or
more missing critical facts.  Keep the domain non-sensitive and the policy
small enough that expected conclusions can be independently enumerated.

Candidate obligations should cover: preserve evidence provenance and polarity;
distinguish ignorance from balanced conflict; prevent two dependent reports
from counting as independent corroboration; apply explicit policy constraints;
expose which premises support a recommendation; and return abstain/Unknown
when a critical threshold or fact is absent.  The benchmark should include a
case where adding evidence retracts or weakens a prior conclusion, exercising
non-monotonic update and downstream invalidation.

\subsubsection*{Expected representations and evidence}

The Plain source should define the decision policy without requiring the
author to write MeTTa.  The IR should represent claims, observations, sources,
dependence, timestamps, weights or evidence measures, policy rules,
contradictions, and explicit unknowns.  MeTTa should own the inspectable
reasoning graph and queries that return conclusions with supporting evidence,
not merely a final label.  Python may supply a grounded evidence aggregation
calculation if its mathematical interpretation is specified independently.

Hypothesis should generate evidence-set permutations, duplicates, dependent
copies, missing facts, and contradictory packets, checking invariance and
abstention properties.  Z3 can test consistency of a bounded crisp policy
fragment.  Lean may prove structural properties such as duplicate-insensitivity
under an explicitly defined aggregation rule.  TLC is likely inapplicable
unless the benchmark adds a workflow in which evidence arrives over time and
actions become irreversible.  No backend should be allowed to manufacture a
ground-truth policy choice that the Plain source leaves unresolved.

\subsubsection*{Mutants and questions the author must settle}

Mutants should: treat missing evidence as false; double-count duplicates;
ignore source dependence; discard negative evidence; overwrite rather than
retain contradiction; threshold before applying a mandatory policy constraint;
lose provenance during aggregation; or convert abstention into the apparent
majority class.  Expected results should explicitly include ignorance,
conflict, supported decision, prohibited decision, and insufficient-policy
states.

The author must choose the evidence semantics, dependence representation,
revision behavior, threshold ownership, explanation requirements, and the
boundary between a mathematical aggregation result and a normative decision.
If these choices are still research questions, the benchmark should expose
them as variants rather than pick one invisibly.

\subsection{Portfolio-level acceptance}

The four cases should share artifact schemas and backend interfaces but should
not be forced into one universal semantic profile.  Portfolio acceptance
should require at least four distinct semantic shapes to traverse the general
pipeline without exact-source conditional branches in the orchestrator.  Each
case should reconstruct its ancestry, kill an agreed preregistered set of
relevant mutants, preserve justified Unknowns, and render tool scopes and
assumptions in the API/UI.  Surviving mutants must be reviewed and classified;
the score itself is not a probability of correctness.

The benchmark-writing phase is complete only when a reviewer could implement
the cases without inventing the missing semantics and could tell, before
running the system, which observations would support, refute, or leave each
obligation unresolved.

\section{Two work plans for the next development phase}

Both plans begin from the merged prototype and its accepted Stage 1--11
vertical.  They are not estimates for rebuilding that infrastructure.  They
cover authoring the realistic benchmarks, making the minimum generalizations
those benchmarks require, integrating the four verticals, and producing
reproducible acceptance evidence.  Calendar estimates assume local tools and
dependencies remain available and no new production deployment, paid compute,
or unrestricted model-provider integration enters scope.

\subsection{Plan S: one coding agent}

The single-agent plan optimizes for semantic continuity and low coordination
cost.  Its principal risk is a self-confirming implementation: one agent may
write the interpretation, lowering, oracle, and test in mutually compatible
ways.  The plan therefore inserts explicit review and mutation checkpoints and
requires the agent to freeze intent artifacts before implementation.

\begin{longtable}{@{}p{0.10\textwidth}p{0.23\textwidth}p{0.49\textwidth}p{0.12\textwidth}@{}}
\toprule
Phase & Deliverable & Work and exit gate & Indicative time \\
\midrule
\endhead
S0 & Baseline and interface map & Reproduce current focused/full suites and all three live examples; map extension seams; record exact branch, commit, tools, and known dirty state. & 1--2 days \\
S1 & Benchmark design pack & Write the four complete Plain tasks, obligation matrices, backend ownership, fixtures, expected Unknowns, and preregistered mutants.  Obtain review before code. & 4--6 days \\
S2 & Shared semantic spine & Generalize exact-source admission, profile dispatch, contract-pattern registry, fixture loading, and evidence projection only as required by the frozen designs.  Preserve old examples. & 3--5 days \\
S3 & Transaction benchmark & Implement the end-to-end reservation/payment slice, focused properties/model, mutants, MeTTa/Python observations, and UI evidence. & 4--6 days \\
S4 & Authorization benchmark & Implement the concurrent revocation slice and its independent oracles/mutants. & 4--6 days \\
S5 & ML benchmark & Implement temporal lineage, numerical runner, statistical procedure, properties, and honest scientific/operational boundaries. & 4--6 days \\
S6 & Epistemic benchmark & Implement evidence/provenance reasoning, conflict/ignorance cases, and abstention-preserving oracles. & 4--6 days \\
S7 & Cross-case integration & Remove case-specific core branches, reconcile schemas, run all cases after every relevant source/contract mutation, and review surviving mutants. & 3--5 days \\
S8 & Acceptance and report & Clean-checkout full suite, pinned backends, trace reconstruction, UI smoke, hygiene and secret checks, independent audit, and final documentation. & 2--4 days \\
\bottomrule
\end{longtable}

This gives roughly 29--46 focused working days, or six to nine calendar weeks
allowing for review and debugging.  The sequence is intentionally benchmark by
benchmark: the agent should finish a thin vertical before expanding the next,
so infrastructure is generalized from observed needs.  If a common abstraction
emerges during S3, it should be validated against at least the existing examples
and the next benchmark design before becoming core architecture.

The agent should maintain one task branch or isolated worktree, commit coherent
increments, and keep a running experiment ledger.  Each phase begins with a
red acceptance test or an immutable design artifact and ends with focused
tests, broader regression, diff hygiene, and a project-record update.  At S1,
S4, S6, and S8, use an independent reviewer---human or separately prompted
frontier model with frozen inputs---to challenge semantics rather than merely
style.

\subsection{Plan P: six parallel coding agents}

The six-agent plan exploits the modular seams already built.  It does not run
the old 11 stages independently six times.  It freezes shared contracts, then
assigns vertical ownership so most files and semantic decisions are disjoint.
One agent owns integration continuously; benchmark agents do not casually edit
the shared composer, store, or public API schemas.

\begin{center}
\small
\begin{tabularx}{0.98\textwidth}{@{}>{\bfseries}p{0.09\textwidth}p{0.23\textwidth}X@{}}
\toprule
Agent & Primary ownership & Responsibilities \\
\midrule
P0 & Architecture/integration & Freeze shared schemas and interfaces; maintain integration branch; review any core change; run continuous contract/regression tests; resolve cross-case ancestry and API/UI consistency. \\
P1 & Transaction case & Author and implement saga contracts, TLC/Hypothesis schedules, local proofs/checks, reference runtime, mutants, and evidence views within its vertical. \\
P2 & Authorization case & Author and implement identity/policy/revocation semantics, concurrency model, policy checks, mutants, and evidence views. \\
P3 & Temporal ML case & Author data/split/metric/uncertainty contract, deterministic local dataset and runner, leakage properties, mutants, and scientific claim boundaries. \\
P4 & Epistemic policy case & Author evidence/conflict/ignorance semantics, MeTTa reasoning graph, aggregation oracle, abstention properties, mutants, and evidence views. \\
P5 & Adversarial/acceptance & Independently review benchmark obligations; build cross-case mutation and stale-artifact attacks; maintain clean-checkout acceptance, reproducibility records, and later documentation/audit. \\
\bottomrule
\end{tabularx}
\end{center}

\subsubsection*{Day 0--2: freeze the collaboration contract}

P0 publishes a versioned benchmark interface pack: canonical artifact fields,
contract-pattern extension interface, backend-result envelope, fixture layout,
G-vector rules, UI projection contract, and ownership map.  P1--P4 draft their
benchmark narratives and obligation matrices in parallel.  P5 critiques them,
looking especially for self-confirming oracles, missing Unknown states, and
claims assigned to inappropriate tools.  At the gate, every requirement has a
stable ID and expected evidence class; unresolved semantics are explicit; the
six agents agree which files are shared and which are vertically owned.

This two-day freeze is the main serial prerequisite.  Implementation may begin
on private fixtures earlier, but no agent may merge a core interface assumption
that has not passed the gate.

\subsubsection*{Day 3--7: four verticals and the shared spine in parallel}

P1--P4 implement the smallest end-to-end slice of their case in independent
worktrees.  Each slice includes source, contracts, one MeTTa observation, one
grounded observation, at least one independent property or formal check, one
deliberately Unknown obligation, and two semantic mutants.  P0 extracts only
the shared generalizations needed to admit those slices and integrates at
least daily.  P5 begins adversarial tests against each accepted interface and
checks that source or contract changes invalidate all appropriate descendants.

Benchmark branches rebase or merge from the integration spine frequently;
they do not merge each other.  A change to a shared schema requires P0 review,
compatibility tests for the old examples, and notification to all owners.
Generated fixtures and golden outputs remain case-owned data rather than
conditional code in the orchestrator.

\subsubsection*{Day 6--10: deepen evidence while integrating continuously}

Once a thin slice is accepted, P1--P4 add the remaining state shapes, backend
tasks, mutants, counterexample presentation, and exact traceability.  P5 runs
cross-case attacks: stale plan reuse, swapped evidence, incorrect backend
attribution, shared-dependency overpromotion, runtime disagreement, timeout,
and semantic-preserving reformatting versus meaning-changing edits.  P0 removes
accidental case coupling, stabilizes API/UI rendering, and runs the full suite
on the integration branch.

The gate is not ``all agents report done.''  It is one integrated commit from
which every old and new case traverses the general pipeline, every obligation
has a server-derived verdict, all preregistered mutants have a recorded
disposition, and justified Unknowns remain visible.

\subsubsection*{Day 10--12: convergence and independent challenge}

Feature work stops except for failures found by acceptance.  P5 owns the
clean-checkout candidate and reproduces the suite from recorded commands.  P0
audits shared code and ancestry.  P1 reviews P2's semantics, P2 reviews P3,
P3 reviews P4, and P4 reviews P1, so no benchmark is accepted solely by its
author.  Reviews compare Plain requirements, contracts, oracles, mutants, and
rendered evidence; they do not merely inspect code style.

Surviving mutants are classified as equivalent, out of scope, oracle gaps, or
implementation defects with evidence.  Any critical oracle gap reopens the
owning vertical.  Tool timeouts or inapplicability remain explicit results and
are not patched into Pass.

\subsubsection*{Day 12--15: release-candidate evidence and documentation}

P5 runs the final clean-checkout acceptance and artifact/hash checks.  P0
prepares the integrated diff and rollback/restart notes.  Benchmark owners
finalize exact case documentation and reproduce one another's commands.
An independent frontier-model or human audit receives frozen source,
contracts, results, and known limitations.  Findings receive explicit
accept/reject/defer decisions.  Merge, deployment, or public exposure remains
a separate authorization decision.

Under these assumptions the realistic target is 10--15 working days, or about
two to three calendar weeks.  Six agents do not provide a sixfold speedup:
the two-day contract freeze, integration latency, final mutation campaign, and
independent audit remain serial.  The gain comes from parallel benchmark
verticals and continuous integration, not from assigning multiple agents to
edit the same core files.

\subsection{Coordination controls for the parallel plan}

The parallel plan should use the following controls:

\begin{itemize}
\item one worktree and task branch per agent, with a declared file-ownership
      map and no direct pushes to the default branch;
\item one versioned interface pack and contract-test suite that all verticals
      consume;
\item P0 approval for shared schemas, canonical serialization, composer,
      artifact-store, and public API changes;
\item daily integration rather than a final ``big bang'' merge;
\item machine-readable fixtures and expected evidence, minimizing merge
      conflicts in central code;
\item focused tests on each branch and the full provider-free/vertical suite
      on the integration branch;
\item preregistered semantic mutants owned or independently reviewed by P5;
\item exact commit/hash references in every handoff, so agents cannot validate
      different versions while believing they share a baseline;
\item stop-the-line treatment for provenance, invalidation, evidence
      attribution, and overpromotion defects; and
\item a final cross-review rotation and clean-checkout reproduction before any
      completion claim.
\end{itemize}

If only three or four agents are available, retain P0 integration and P5
adversarial independence, then combine benchmark ownership by semantic
affinity.  If more than six are available, add reviewers, fixture/oracle
authors, and formal-backend specialists rather than splitting the shared core
among more simultaneous editors.  Coordination quality, not raw agent count,
sets the useful parallelism ceiling.

\end{document}
